\documentclass[11pt,a4paper]{article}

% ===================================================
% Project metadata
% ===================================================
\newcommand{\documentTitle}{Advanced RV32I Core with Cache Subsystem}
\newcommand{\coverTitle}{Advanced RV32I Core with\\Cache Subsystem}
\newcommand{\documentVersion}{v0.1}
\newcommand{\authorName}{Nguyen To Quoc Viet}

% ===================================================
% Packages
% ===================================================
\usepackage[a4paper, top=2.5cm, bottom=2.5cm, left=2cm, right=2cm]{geometry}
\usepackage[utf8]{inputenc}
\usepackage[T1]{fontenc}
\usepackage[vietnamese,english]{babel}
\usepackage{vntex}

\usepackage{amsmath, amssymb, mathtools}
\usepackage{graphicx}
\usepackage{xcolor}
\usepackage{tikz}
\usetikzlibrary{positioning, arrows.meta, shapes.geometric, fit, calc}
\usepackage{float}

\usepackage{indentfirst}
\usepackage{setspace}
\usepackage{enumitem}
\usepackage{array, booktabs}
\usepackage{titlesec}
\usepackage{fancyhdr}

\usepackage{listings}
\definecolor{codegreen}{rgb}{0,0.6,0}
\definecolor{codegray}{rgb}{0.5,0.5,0.5}
\definecolor{codepurple}{rgb}{0.58,0,0.82}
\definecolor{backcolour}{rgb}{0.96,0.96,0.94}

\lstdefinelanguage{SystemVerilog}{
  keywords={module, endmodule, input, output, logic, always_ff, always_comb,
            posedge, negedge, begin, end, if, else, case, endcase, assign,
            localparam, parameter, import, package, endpackage, typedef, enum,
            for, int, always, wire, reg},
  keywordstyle=\color{blue!70!black}\bfseries,
  comment=[l]{//},
  morecomment=[s]{/*}{*/},
  commentstyle=\color{codegreen}\itshape,
  stringstyle=\color{codepurple},
  sensitive=true
}

\lstset{
  language=SystemVerilog,
  backgroundcolor=\color{backcolour},
  basicstyle=\ttfamily\small,
  breaklines=true,
  captionpos=b,
  keepspaces=true,
  numbers=left,
  numberstyle=\tiny\color{codegray},
  numbersep=5pt,
  tabsize=2,
  frame=single,
  framerule=0.3pt
}

\usepackage[hidelinks]{hyperref}

% ===================================================
% Header / footer
% ===================================================
\setlength{\headheight}{28pt}
\pagestyle{fancy}
\fancyhead{}
\fancyhead[L]{
  \begin{tabular}{l}
    \ttfamily\bfseries \documentTitle
  \end{tabular}
}
\fancyhead[R]{\ttfamily\bfseries \documentVersion}

\fancyfoot{}
\fancyfoot[L]{\scriptsize\ttfamily \authorName}
\fancyfoot[R]{\scriptsize\ttfamily Page \thepage}
\renewcommand{\headrulewidth}{0.3pt}
\renewcommand{\footrulewidth}{0.3pt}

% Paragraph title format
\titleformat{\paragraph}
{\normalfont\normalsize\bfseries}{\theparagraph}{1em}{}
\titlespacing*{\paragraph}
{0pt}{3.25ex plus 1ex minus .2ex}{1.5ex plus .2ex}

% ===================================================
% Document
% ===================================================
\begin{document}

\begin{titlepage}
\centering
\vspace*{2.5cm}

\rule{0.8\linewidth}{0.6pt}
\vspace{0.6cm}

{\Huge\bfseries \coverTitle \par}
\vspace{0.4cm}
{\Large\ttfamily \documentVersion \par}

\vspace{0.6cm}
\rule{0.8\linewidth}{0.6pt}

\vfill

{\large\bfseries Author\par}
\vspace{0.3cm}
{\Large \authorName \par}

\end{titlepage}

\newpage
\renewcommand{\contentsname}{Table of Contents}
\tableofcontents
\pagebreak

\section{Scope}

\subsection{Core Scope}

The core documented in this report is the currently implemented RV32I
five-stage in-order processor. It uses a single-issue pipeline with the
classic IF, ID, EX, MEM, and WB stages. The implemented instruction set is
the RV32I base integer ISA, including integer arithmetic, logical
operations, shifts, comparisons, load/store instructions, conditional
branches, JAL, JALR, LUI, and AUIPC.

The core includes dynamic branch prediction using a 1024-entry Branch
History Table and Branch Target Buffer. Branch and jump instructions are
resolved in the EX stage, and mispredictions redirect the fetch PC through
a registered feedback path. Data hazards are handled by forwarding from
later pipeline stages and by inserting a stall for load-use hazards.

\subsection{Cache Subsystem Scope}

The cache subsystem connects the processor core to an external AXI4 memory
interface. It contains a 4 KB direct-mapped instruction cache, a 4 KB
two-way set-associative data cache, a 4-entry write buffer, and a bus
arbiter shared by instruction fetches, data-cache refills, and write-buffer
drains. Both caches use 16-byte cache lines and a 32-bit data path.

The write buffer decouples stores from the external memory path and also
provides store-to-load forwarding for pending writes. Cache misses are
serviced through the AXI4 master interface.

\subsection{Boundary}

This document covers only the implemented RV32I five-stage baseline and
its cache subsystem. RV32M, floating-point, vector/IME extensions, CSR,
exception handling, interrupts, privilege levels, MMU, virtual memory, and
full cache coherency are outside the implemented baseline.

\section{Source Repository}

The source code of this project is maintained in the following GitHub repository:

\begin{center}
\href{https://github.com/NguyenToQuocViet/Advanced-RISC-V-32-bit-SoC}
{\texttt{github.com/NguyenToQuocViet/Advanced-RISC-V-32-bit-SoC}}
\end{center}

The architecture described in this document corresponds to the integrated
five-stage RV32I baseline in the RTL source tree, especially the modules under
\texttt{rtl/cpu/core/} and \texttt{rtl/cpu/cache/}.

\clearpage
\section{Top-Level Architecture}

\begin{figure}[H]
\centering
\resizebox{0.78\textwidth}{!}{%
\begin{tikzpicture}[
    socbox/.style={
        draw,
        thick,
        dashed,
        rounded corners,
        minimum width=8.2cm,
        minimum height=8.0cm,
        inner sep=0pt
    },
    block/.style={
        draw,
        thick,
        rounded corners,
        minimum width=5.4cm,
        minimum height=1.25cm,
        align=center,
        font=\small\ttfamily
    },
    extblock/.style={
        draw,
        thick,
        rounded corners,
        minimum width=5.0cm,
        minimum height=1.25cm,
        align=center,
        font=\small\ttfamily
    },
    bus/.style={-{Latex[length=2.2mm]}, thick},
    ctrl/.style={-{Latex[length=2.2mm]}, thick, dashed},
    buslabel/.style={font=\scriptsize\ttfamily, fill=white, inner sep=2pt}
]

\node[socbox] (soc) at (0,0) {};
\node[font=\small\ttfamily\bfseries, anchor=west] at ([xshift=0.25cm,yshift=-0.3cm]soc.north west) {riscv\_soc};

\node[block] (core) at (0,1.8) {riscv\_core};
\node[block] (cache) at (0,-1.3) {cache\_subsystem};
\node[extblock] (mem) at (0,-5.6) {behavioral memory model};

\draw[bus] ([xshift=-1.25cm]core.south) -- node[buslabel, left] {IF channel} ([xshift=-1.25cm]cache.north);
\draw[bus] ([xshift=1.25cm]core.south) -- node[buslabel, right] {MEM channel} ([xshift=1.25cm]cache.north);
\draw[bus] (cache.south) -- node[buslabel, right] {AXI4 master} (mem.north);

\end{tikzpicture}%
}
\caption{Top-level organization of the implemented RV32I core and cache subsystem.}
\end{figure}

At the top level, the core-cache boundary is intentionally split into two
independent channels. The IF channel carries instruction-fetch requests from
the core to the cache subsystem, while the MEM channel carries load and store
requests. This separation mirrors the pipeline structure: instruction traffic
originates from the IF stage, whereas data-memory traffic originates from the
MEM stage.

The IF channel carries the requested program counter and the returned
instruction. The signal \texttt{if\_req} indicates that the core requests the
instruction at \texttt{if\_pc}. The cache subsystem returns
\texttt{if\_instr}, \texttt{if\_icache\_valid}, and
\texttt{if\_icache\_ready}, which determine whether the fetch stage can
consume the instruction and advance the program counter.

The MEM channel carries the effective address, request type, store data, and
byte write strobes. Loads and stores share \texttt{mem\_addr} and
\texttt{mem\_req}. Stores additionally use \texttt{mem\_we},
\texttt{mem\_wdata}, and \texttt{mem\_wstrb}. The cache subsystem returns
\texttt{mem\_rdata}, \texttt{mem\_dcache\_valid}, and
\texttt{mem\_dcache\_ready}.

The cache subsystem is the only block that drives the external AXI4 master
interface. At this level, AXI4 is treated as five grouped channels: read
address, read data, write address, write data, and write response. Detailed
AXI4 signal behavior is described with the cache subsystem rather than at the
top-level boundary.

\begin{table}[H]
\centering
\small
\begin{tabular}{@{}p{2.2cm}p{2.5cm}p{7.0cm}p{2.3cm}@{}}
\toprule
\textbf{Interface} & \textbf{Direction} & \textbf{Main signals} & \textbf{Purpose} \\
\midrule
IF channel & Core $\leftrightarrow$ cache &
\texttt{if\_pc}, \texttt{if\_req}, \texttt{if\_instr},
\texttt{if\_icache\_ready}, \texttt{if\_icache\_valid} &
Instruction fetch \\
\addlinespace
MEM channel & Core $\leftrightarrow$ cache &
\texttt{mem\_addr}, \texttt{mem\_req}, \texttt{mem\_we},
\texttt{mem\_wdata}, \texttt{mem\_wstrb}, \texttt{mem\_rdata},
\texttt{mem\_dcache\_ready}, \texttt{mem\_dcache\_valid} &
Load/store access \\
\addlinespace
AXI4 master & Cache $\leftrightarrow$ memory &
AR, R, AW, W, and B channels & External memory transactions \\
\bottomrule
\end{tabular}
\caption{Top-level interfaces between the core, cache subsystem, and memory.}
\end{table}

\clearpage
\section{Five-Stage Pipeline}

\section{Core Modules}

\section{Cache Subsystem}

\section{Verification Surface}

\section{Known Limitations}

\end{document}
